% page 331 du fac-simile (image p-330 du scan)
% bloc 170.2 x 237.7 mm ; marge gauche 31.8 mm
\pgc{10.160mm}{0.000mm}{
\l{\cel{56}323}
\l{}
\l{\sou{kuraro}. Veneno ek planti, aden qua la Amerindiani di la}
\l{\cel{3}sudo-regiono di Amerika trempas sua flechi. - DEFIS.}
\l{}
\l{\sou{kuraso}. I. Armo defensala ek ledro, metalo, e c., qua envol-}
\l{\cel{3}vas e protektas la pektoro e la dorso. - II. La squamaro}
\l{\cel{3}qua envolvas ula fishi. - III. Kovrilo ek plaki de fero}
\l{\cel{3}o de stalo qua protektas la kareno di ula milito-navi. - DEFIRS.}
\l{}
\l{\sou{kuratar}. (\sou{yurosienco}). Komisesar, segun la lego, pri defensor la}
\l{\cel{3}havajo e la interesti sive di minoro, sive di majoro qua}
\l{\cel{3}esas deklarita nekapabla administrar sua heredajo. - DEFIRS.}
\l{}
\l{\sou{kurbaturar}. (\sou{trans.)}\cel{2}(\sou{aludante fatigo-kauzo ecesanta, o morbo)}.}
\l{\cel{3}Igar extreme nekapabla pri esforco, pro la doloro specala}
\l{\cel{3}qua afektas la muskuli pos granda esforco. - F.}
\l{}
\l{\sou{kurear}. (\sou{netrans.}) (\sou{aludante chaso kun hundi}). Abandonar a la}
\l{\cel{3}chaso-hundi parto de la bestio quan li kaptis. - F.}
\l{}
\l{\sou{kuretar}. (\sou{trans.})\cel{2}Deprenar la kalkoli de la veziko, lor la}
\l{\cel{3}operaco \textquotedbl{}litotomio\textquotedbl{}, per instrumento di qua un extremajo}
\l{\cel{3}esas kulieratra, e qua uzesas por deprenar la materii qui}
\l{\cel{3}normale ne esas en ta kavajo. - F.}
\l{}
\l{\sou{kurio}. (\sou{eklezio katol}.)Ensemblo ek la tribunali, ek la organismi}
\l{\cel{3}administrala, di la guverno-sistemo papala. - DEFIRS.}
\l{}
\l{\sou{kuriero}. I. Persono qua vehas rapide, per kavalo, por arivar}
\l{\cel{3}ante veturo o por preparar la relayi. - II. Portisto di de-}
\l{\cel{3}peshi qua cirkulas pede, kavale o bicikle. - III. Portisto di}
\l{\cel{3}depeshi qua cirkulas per veturo ordinara, automobila, o per}
\l{\cel{3}veturo postala. - DEFIRS.}
\l{}
\l{\sou{kurioza}. Qua deziras vidar, konocar, saveskar. - DEFIRS.}
\l{}
\l{\sou{kurkuliono}. (\sou{zool.}) Insekto koleoptera, di qua un speco rodas}
\l{\cel{3}frumento. - L.\cel{1}\sou{curculio}. - ISL.}
\l{}
\l{\sou{kurkumo}. (\sou{bot.}) Genero di planto herbacea, perena, quan onu ren-}
\l{\cel{3}kontras en Azia, Afrika ed Amerika; la radiko di un speco di}
\l{\cel{3}olu (la safrano di India) uzesas en la tinteyi, pro olua}
\l{\cel{3}kolorizivo flava. - L.\cel{1}\sou{curcuma longa}.\cel{2}- DEFIRS.}
\l{}
\l{kurlio. (\sou{zool.}) Ucelo ek la familio \textquotedbl{}stiltieri\textquotedbl{} longa-beka.}
\l{\cel{3}L.\cel{1}\sou{numenius}. - EFIS.}
\l{}
\l{kurso. I. Serio de lecioni qui formacas ensemblo docala. -}
\l{\cel{3}II. Preco di kompro o di vendo, pri kambiajo. - III. Cirkulo}
\l{\cel{3}regulala di kambiajo. - DEFIRS.}
\l{}
\l{kursiva. (\sou{pri skribo)}. Qua konsistas ek literi trasebla fluante,}
\l{\cel{3}ed inklinata de dextre a sinistre. Dicesas anke pri la imprimo-}
\l{\cel{3}tipo qua inklinesas same kam la skriburo ordinara. - DEFIS.}
}
